\begin{lemma}[Minty's characterization of the proximal point]
Let $E$ be a finite-dimensional real inner-product space, let $\psi:E\to\R$ be
convex on $E$ and let $y,x_0\in E$. Then the following two statements are
equivalent:
\begin{enumerate}
\item $x_0$ minimizes $x\mapsto\psi(x)+|x-y|^2/2$ over $E$, that is
\[
  \psi(x_0)+\frac{|x_0-y|^2}{2}\le\psi(x)+\frac{|x-y|^2}{2}\qquad(x\in E);
\]
\item $y-x_0$ is a subgradient of $\psi$ at $x_0$, that is
\[
  \psi(x_0)+\langle y-x_0,x-x_0\rangle\le\psi(x)\qquad(x\in E).
\]
\end{enumerate}
\end{lemma}
\begin{proof}
Throughout we use the identity, valid for all $x$,
\begin{equation}\label{eq:mps-expand}
  \frac{|x-y|^2}{2}
   =\frac{|(x-x_0)+(x_0-y)|^2}{2}
   =\frac{|x-x_0|^2}{2}+\langle x-x_0,x_0-y\rangle+\frac{|x_0-y|^2}{2},
\end{equation}
which is the expansion of the square of a sum in an inner-product space.

\emph{(2) implies (1).}
Assume (2) and let $x\in E$. Adding $|x-y|^2/2$ to both sides of (2) and using
\eqref{eq:mps-expand},
\[
  \psi(x)+\frac{|x-y|^2}{2}
   \ge\psi(x_0)+\langle y-x_0,x-x_0\rangle
      +\frac{|x-x_0|^2}{2}+\langle x-x_0,x_0-y\rangle+\frac{|x_0-y|^2}{2}.
\]
The two inner products cancel, since
$\langle x-x_0,x_0-y\rangle=-\langle y-x_0,x-x_0\rangle$. Hence
\[
  \psi(x)+\frac{|x-y|^2}{2}
   \ge\psi(x_0)+\frac{|x_0-y|^2}{2}+\frac{|x-x_0|^2}{2}
   \ge\psi(x_0)+\frac{|x_0-y|^2}{2},
\]
which is (1).

\emph{(1) implies (2).}
Assume (1), let $x\in E$ and let $t\in(0,1]$. Put $x_t=x_0+t(x-x_0)$. Applying
(1) at the point $x_t$,
\begin{equation}\label{eq:mps-at-xt}
  \psi(x_0)+\frac{|x_0-y|^2}{2}\le\psi(x_t)+\frac{|x_t-y|^2}{2}.
\end{equation}
Convexity of $\psi$ on $E$ with the weights $1-t$ and $t$ gives
\[
  \psi(x_t)\le(1-t)\psi(x_0)+t\psi(x).
\]
Also $x_t-y=(x_0-y)+t(x-x_0)$, so
\[
  \frac{|x_t-y|^2}{2}
   =\frac{|x_0-y|^2}{2}+t\langle x_0-y,x-x_0\rangle+\frac{t^2|x-x_0|^2}{2}.
\]
Substituting both into \eqref{eq:mps-at-xt} and cancelling
$|x_0-y|^2/2$ from the two sides,
\[
  \psi(x_0)\le(1-t)\psi(x_0)+t\psi(x)
    +t\langle x_0-y,x-x_0\rangle+\frac{t^2|x-x_0|^2}{2},
\]
that is, after subtracting $(1-t)\psi(x_0)$,
\[
  t\,\psi(x_0)\le t\,\psi(x)+t\langle x_0-y,x-x_0\rangle
     +\frac{t^2|x-x_0|^2}{2}.
\]
Dividing by $t>0$,
\[
  \psi(x_0)\le\psi(x)+\langle x_0-y,x-x_0\rangle+\frac{t\,|x-x_0|^2}{2}.
\]
The left-hand side and the first two terms on the right do not depend on $t$, so
letting $t\downarrow0$ gives
$\psi(x_0)\le\psi(x)+\langle x_0-y,x-x_0\rangle$, i.e.
$\psi(x_0)+\langle y-x_0,x-x_0\rangle\le\psi(x)$. As $x$ was arbitrary this is
(2).
\end{proof}
